% 表紙と目次（全単元版のみ）
\thispagestyle{cover}
\vspace*{18mm}
{\sffamily\color{pone}\large 経営情報 2026　AI利用ガイド}\par
\vspace{3mm}
{\sffamily\bfseries\fontsize{30pt}{38pt}\selectfont 復習教材}\par
\vspace{2mm}
{\sffamily\fontsize{20pt}{26pt}\selectfont 問いを立て、調べ、読み、磨く}\par
\vspace{4mm}
{\color{inktwo}\rule{\linewidth}{0.6pt}}\par
\vspace{4mm}
{\sffamily\small 学籍番号\makebox[40mm]{\hrulefill}\qquad 氏名\makebox[60mm]{\hrulefill}}\par
\vspace{10mm}

\begin{tcolorbox}[enhanced,colback=papertwo,colframe=papertwo,boxrule=0pt,arc=1.5mm,
    left=4mm,right=4mm,top=3mm,bottom=3mm,fontupper=\small]
{\sffamily\bfseries このワークシートの使い方}\par\vspace{1mm}
\begin{itemize}[label=・]
  \item AI利用ガイドの「復習教材」に対応する、5つの単元のワークシートです。単元ごとに1枚（2〜3ページ）で、復習1から順に進めると、一つの問いが確定するまでの流れをたどれます。
  \item 「■ ワーク」は手を動かす課題、「仕上げ」はワークの後に取り組む課題です。目安の時間は単元の見出しにあります。
  \item AIの出力は記入欄にそのまま写さず、要点を自分の言葉で書いてください。
  \item 各単元の末尾の「AI活用記録」は、提出物に添付するAI活用記録の練習になります（書き方はAI利用ガイド）。
  \item 使用するAIは大学が提供する Copilot です。個人情報や企業などの非公開の情報は入力しません。
  \item 記入例を別に用意しています。自分で取り組んでから見比べてください。
\end{itemize}
\end{tcolorbox}

\vspace{6mm}
{\sffamily\bfseries 目次}\par\vspace{1mm}
{\small\makeatletter\@starttoc{toc}\makeatother}
